% Copy the two PDFs into your LaTeX project's figures/ directory.
% Preamble: \usepackage{graphicx}
% Use figure* for full-width figures in the ACL two-column layout.
% These are figure snippets, not a standalone LaTeX document.

\begin{figure*}[t]
  \centering
  \includegraphics[width=\textwidth]{figures/tweetqa_final_architecture.pdf}
  \caption{Final prediction system. (a) The from-scratch TweetBiDAF span
  predictor combines word, character, and token-type features with a shared
  BiLSTM encoder, bidirectional question--context attention, and gated
  self-attention. (b) Three independently initialized models trained on the
  same training partition produce start and end distributions, which are
  averaged separately. The ten highest-scoring valid spans under
  $\log \bar p_s(i)+\log \bar p_e(j)$ are represented by 30 features and
  scored by a fixed OOF-trained LambdaMART reranker. The highest-ranked
  span is extracted from the original tweet. Gold answers are not used at
  prediction time.}
  \label{fig:final-system}
\end{figure*}

\begin{figure*}[t]
  \centering
  \includegraphics[width=\textwidth]{figures/tweetqa_oof_training.pdf}
  \caption{Out-of-fold (OOF) candidate generation for reranker training.
  The training partition is divided into three tweet-grouped folds.
  Each fold model is fitted and selected using only the other two folds,
  with a separate internal validation subset for early stopping.
  Predictions on the held-out fold provide candidate features; gold
  answers in that fold provide relevance labels. Records from all folds
  are pooled to train one reranker. Outer validation and test data remain
  separate. These fold models are distinct from the three full-training-partition
  prediction models in Figure~\ref{fig:final-system}.}
  \label{fig:oof-training}
\end{figure*}
